\section{The locus of Hodge classes}\label{app:cdk}

Every deformation argument about Hodge classes rests on one theorem: the locus
in a family where a given class stays a Hodge class is an algebraic
subvariety, not merely an analytic one. This is the theorem of Cattani,
Deligne and Kaplan \cite{CDK95}. The present appendix states it, proves the
part of it that is elementary, isolates the two analytic inputs that the rest
needs, computes the example that matters here, and then says precisely why the
theorem is not a route to the Hodge conjecture. The last point is the one that
bears on \Cref{sec:construction}, and it is the reason the appendix is here
rather than in a reference.

\subsection{Variations of Hodge structure}

\begin{definition}[Variation of Hodge structure]\label{def:vhs}
Let $S$ be a smooth connected complex quasi-projective variety. A
\emph{polarised variation of $\QQ$-Hodge structure} of weight $w$ on $S$
consists of
\begin{enumerate}[label=\textup{(\alph*)},nosep]
\item a local system $\mathbb{V}$ of finite dimensional $\QQ$-vector spaces on
  $S$, with associated holomorphic bundle
  $\cV=\mathbb{V}\otimes_{\QQ}\cO_{S}$ carrying the flat connection $\nabla$
  whose flat sections are $\mathbb{V}$;
\item a decreasing filtration $\cV=F^{0}\supseteq F^{1}\supseteq\cdots$ by
  holomorphic subbundles, inducing on each fibre $\mathbb{V}_{s}$ a
  $\QQ$-Hodge structure of weight $w$;
\item a flat nondegenerate $(-1)^{w}$-symmetric pairing
  $Q\colon\mathbb{V}\otimes\mathbb{V}\to\QQ_{S}$ polarising each fibre;
\end{enumerate}
subject to Griffiths transversality
\begin{equation}\label{eq:transversality}
  \nabla F^{p}\subseteq F^{p-1}\otimes\Omega^{1}_{S}.
\end{equation}
\end{definition}

The example to keep in mind is the geometric one: for $f\colon\cX\to S$ smooth
projective, $\mathbb{V}=R^{w}f_{*}\QQ$ with its Gauss--Manin connection is such
a variation, and \eqref{eq:transversality} is Griffiths' theorem
\cite[Ch.~10]{Voi02}. Our other example is the Weil family of
\Cref{prop:weildomain} below, which is not defined by a family of varieties in
the first instance but by a period domain, and which carries a variation all
the same.

\begin{notation}\label{not:totalspace}
Write $\cV^{\mathrm{an}}$ for the total space of $\cV$, a complex manifold
fibred over $S$ with fibre $\mathbb{V}_{s}\otimes\CC$, and
$\mathbb{V}^{\mathrm{tot}}\subset\cV^{\mathrm{an}}$ for the subset of pairs
$(s,\alpha)$ with $\alpha\in\mathbb{V}_{s}$, a covering space of $S$ with
countable fibres.
\end{notation}

\subsection{The locus of Hodge classes}

\begin{definition}\label{def:hodgelocus}
Let $\mathbb{V}$ be a polarised variation of weight $2p$ on $S$. Its
\emph{locus of Hodge classes} is
\[
  \Hdg^{p}(\mathbb{V})
  =\bigl\{(s,\alpha)\in\mathbb{V}^{\mathrm{tot}}\;:\;
     \alpha\in F^{p}\cV_{s}\bigr\}.
\]
For $N\in\QQ$ write $\Hdg^{p}_{N}$ for the part with $Q(\alpha,\alpha)=N$.
\end{definition}

The condition $\alpha\in F^{p}$ for a rational $\alpha$ of weight $2p$ is
exactly the condition that $\alpha$ be of type $(p,p)$, by conjugation
symmetry. The following two lemmas are elementary and give everything about
$\Hdg^{p}$ except algebraicity.

\begin{lemma}[Analyticity]\label{lem:hodgeanalytic}
$\Hdg^{p}(\mathbb{V})$ is a closed analytic subset of
$\mathbb{V}^{\mathrm{tot}}$, and the projection
$\pi\colon\Hdg^{p}(\mathbb{V})\to S$ is locally biholomorphic onto its image on
each connected component. In particular each component is a closed analytic
subvariety of an open subset of $S$, and there are countably many components.
\end{lemma}

\begin{proof}
Work on a simply connected open $U\subseteq S$ over which $\mathbb{V}$ is
trivialised, $\mathbb{V}|_{U}\cong V_{0}\times U$ for a fixed $\QQ$-vector
space $V_{0}$. Then $\mathbb{V}^{\mathrm{tot}}|_{U}$ is the disjoint union over
$\alpha\in V_{0}$ of copies of $U$. Fix $\alpha$. The quotient
$\cV/F^{p}$ is a holomorphic vector bundle on $U$ and the image
$\bar\alpha(s)$ of the flat section $\alpha$ in $(\cV/F^{p})_{s}$ is a
holomorphic section of it. The locus in question is the zero locus of
$\bar\alpha$, a closed analytic subset of $U$. Taking the union over the
countable set $V_{0}$ gives the first statement, and shows that on the sheet
indexed by $\alpha$ the projection $\pi$ is the inclusion of that zero locus
into $U$, hence injective. Countability of the set of components follows
because $V_{0}$ is countable and each zero locus has countably many
components.
\end{proof}

\begin{lemma}[Finiteness at bounded norm]\label{lem:boundednorm}
Fix $s\in S$ and $N\in\QQ$. The set of $\alpha\in\mathbb{V}_{s}$ with
$\alpha$ of type $(p,p)$, $Q(\alpha,\alpha)=N$, and $\alpha$ in a fixed
lattice $\mathbb{V}_{s,\ZZ}$, is finite. The number of such $\alpha$ is locally
constant on each component of $\Hdg^{p}_{N}$.
\end{lemma}

\begin{proof}
Decompose $\alpha$ into its primitive parts for the Lefschetz decomposition
attached to the polarisation. On the primitive part of type $(p,p)$ in weight
$2p$ the Hodge--Riemann bilinear relations read
$(-1)^{p}\,Q(\alpha,\bbar{\alpha})>0$ for $\alpha\ne0$, so the form
$(-1)^{p}Q$ is positive definite on the real span of the primitive real
$(p,p)$-classes; the same holds on each Lefschetz summand with an alternating
sign, and the resulting form $\tilde Q$ obtained by flipping signs summand by
summand is positive definite on the whole space of real $(p,p)$-classes. The
conditions $Q(\alpha,\alpha)=N$ and $\alpha$ of type $(p,p)$ therefore confine
$\alpha$ to a compact set for $\tilde Q$, and a compact set meets a lattice in
finitely many points. The count is locally constant because the lattice and $Q$
are flat and the Hodge filtration varies continuously, so no Hodge class can
appear or disappear without its $Q$-norm changing.
\end{proof}

\begin{theorem}[Cattani, Deligne and Kaplan \cite{CDK95}]\label{thm:cdk}%
\footnote{The paper is short and the proof is not. It draws on the whole asymptotic theory of degenerating Hodge structures built in the two decades before it, and the two inputs isolated below are the load-bearing ones. Deligne had earlier shown that Hodge classes on abelian varieties are absolute by a different route; the theorem here is what carries the relevant consequences beyond abelian varieties.}
Let $S$ be a smooth complex quasi-projective variety and $\mathbb{V}$ a
polarised variation of $\ZZ$-Hodge structure of weight $2p$ on $S$. Then each
connected component of $\Hdg^{p}(\mathbb{V})$ is a closed algebraic subvariety
of $\mathbb{V}^{\mathrm{tot}}$, and its image in $S$ is a closed algebraic
subvariety of $S$.
\end{theorem}

We do not reproduce the proof. What we do record is exactly where the
difficulty lies and which two theorems remove it, because the shape of the
argument is what a reader needs in order to judge whether a deformation
argument in this area is legitimate.

By \Cref{lem:hodgeanalytic} a component $Z$ of $\Hdg^{p}$ is already a closed
analytic subvariety of an open subset of $S$, and by \Cref{lem:boundednorm} it
carries a locally constant $Q$-norm and is one sheet of a finite cover of its
image. Choose a smooth projective compactification $\bar S\supseteq S$ with
$\bar S\setminus S$ a normal crossing divisor. By Chow's theorem it is enough
to show that $Z$ extends to an analytic subvariety of $\bar S$, that is,
that $Z$ does not oscillate wildly as one approaches the boundary. That is an
asymptotic question about the Hodge filtration, and the two theorems that
answer it are:

\begin{enumerate}[label=\textup{(\Alph*)},leftmargin=2.2em]
\item \emph{The nilpotent orbit theorem} of Schmid \cite{Sch73}: near a
  boundary point the period map is approximated, to exponentially small error
  in the natural coordinates, by the orbit of a nilpotent Lie algebra element
  built from the local monodromy logarithms.
\item \emph{The $\SL_{2}$-orbit theorem} in several variables, of Cattani,
  Kaplan and Schmid \cite{CKS86}: that nilpotent orbit is in turn approximated
  by an orbit of a fixed real reductive group, with uniform estimates in the
  several-variable setting.
\end{enumerate}

Together these give a two-sided estimate for the Hodge norm of a flat section
near the boundary, of the shape $\|\alpha\|_{s}\sim\prod_{j}(\log|t_{j}|^{-1})^{m_{j}}$
with integers $m_{j}$ determined by the weight filtration. A flat section that
is Hodge on $Z$ has $\tilde Q$-norm bounded by \Cref{lem:boundednorm}, and the
estimate converts that bound into a statement that the possible boundary
behaviours form a finite list. Algebraicity then follows from the analytic
extension across the boundary together with Chow's theorem in $\bar S$.

\begin{remark}\label{rem:cdkinputs}
Both (A) and (B) are theorems about degenerating Hodge structures and neither
mentions algebraic cycles. This is worth saying plainly: the algebraicity in
\Cref{thm:cdk} is algebraicity of a locus of periods, obtained from analysis,
and it carries no information about algebraic cycles whatever. The point is
developed in \Cref{ssec:cdknot}.
\end{remark}

\subsection{The Weil family}

The relevant example for this paper is the family of complex structures in
\Cref{setup:model}.

\begin{proposition}[The Weil period domain]\label{prop:weildomain}
Fix $K$, $n$, the $K$-space $V=K^{2n}$, and the hermitian form $H$ of
signature $(n,n)$ from \Cref{lem:hermform}. The set of Hodge structures of
weight one on $V$ that are polarised by $Q$, commute with the $K$-action, and
are of $(K,-1,n)$-Weil type is in bijection with
\[
  \cD_{n,n}
  =\bigl\{P\subset V_{+}\;:\;\dim_{\CC}P=n,\ H|_{P}\ \text{negative definite}\bigr\},
\]
an open subset of the Grassmannian $\mathrm{Gr}(n,2n)$ and a bounded symmetric
domain isomorphic to $\mathrm{U}(n,n)/\bigl(\mathrm{U}(n)\times\mathrm{U}(n)\bigr)$.
In particular
\[
  \dim_{\CC}\cD_{n,n}=n^{2},
\]
against $\dim\cA_{2n}=n(2n+1)$ for the moduli of principally polarised abelian
$2n$-folds.
\end{proposition}

\begin{proof}
A Hodge structure of weight one on $V$ is a subspace $V^{1,0}\subset V_{\CC}$
with $V_{\CC}=V^{1,0}\oplus\overline{V^{1,0}}$. Commuting with the $K$-action
means that $V^{1,0}$ is stable under $S$, hence splits as
$V^{1,0}=P\oplus P'$ with $P=V^{1,0}\cap V_{+}$ and $P'=V^{1,0}\cap V_{-}$.
Conjugation interchanges $V_{+}$ and $V_{-}$, so
$\overline{V^{1,0}}=\bbar{P}\oplus\bbar{P'}$ with $\bbar{P}\subset V_{-}$, and
the direct sum condition becomes $V_{-}=P'\oplus\bbar{P}$, which determines
$P'$ up to the choice of a complement; the Weil condition
$\dim P=\dim P'=n$ of \Cref{lem:balanced} together with the polarisation
condition pins it down as $P'=\bbar{P}^{\perp_{H}}$. So the structure is
determined by $P$.

For the positivity, extend $H$ to the $\CC$-bilinear pairing between $V_{+}$
and $V_{-}$; the computation \eqref{eq:Qplusminus} in the proof of
\Cref{thm:model} says that the second Riemann relation holds exactly when
$i\,Q$ is positive on $V^{1,0}$, which on the $V_{+}$ part is the condition
that $H$ be negative definite on $P$, the sign being fixed by the convention
in \Cref{lem:hermform}. The set of $n$-dimensional negative definite subspaces
of a hermitian space of signature $(n,n)$ is the standard realisation of the
symmetric domain of $\mathrm{U}(n,n)$: the unitary group acts transitively on
it and the stabiliser of a point is $\mathrm{U}(n)\times\mathrm{U}(n)$, the
unitary groups of $P$ and of $P^{\perp}$. It is open in $\mathrm{Gr}(n,2n)$
because negative definiteness is an open condition, and
$\dim\mathrm{Gr}(n,2n)=n(2n-n)=n^{2}$.

The comparison dimension is $\dim\cA_{g}=g(g+1)/2$ with $g=2n$, which is
$n(2n+1)$.
\end{proof}

\begin{corollary}[The Weil locus is algebraic]\label{cor:weilalgebraic}
Inside $\cA_{2n}$ the image of $\cD_{n,n}$, that is the locus of principally
polarised abelian $2n$-folds admitting an action of an order of $K$ of
$(K,-1,n)$-Weil type, is a closed algebraic subvariety of dimension $n^{2}$.
\end{corollary}

\begin{proof}
The condition of admitting such an action is the condition that a certain
class in $\End(H^{1})$, namely the graph of $\iota(\sqrt{-d})$, be a Hodge
class, together with the balanced signature, which is locally constant. So the
locus is a component of a locus of Hodge classes for the variation
$\End(R^{1}f_{*}\QQ)$ over $\cA_{2n}$ with its level structure, and
\Cref{thm:cdk} applies. The dimension is \Cref{prop:weildomain}.
\end{proof}

\subsection{What the theorem does not give}\label{ssec:cdknot}

\begin{remark}[Algebraic locus, not algebraic class]\label{rem:notalgebraic}
\Cref{thm:cdk} says that the set of parameters where $\alpha$ is Hodge is cut
out by polynomial equations. It says nothing about whether $\alpha$ is the
class of a cycle at any one of those parameters. The two statements are
independent in both directions: the Hodge locus of a class that happens to be
algebraic everywhere is still only known to be algebraic by \Cref{thm:cdk},
and conversely knowing the locus is algebraic supplies no cycle. It is easy to
slide from one to the other in words, and \Cref{cor:weilalgebraic} is exactly
the kind of statement that invites the slide: the Weil locus is an algebraic
subvariety, it is even a Shimura subvariety, and none of that produces the
class $\omega_{1}$ of \eqref{eq:omegagens}.
\end{remark}

\begin{remark}[What would be needed: the variational conjecture]%
\label{rem:variational}
The property that would turn \Cref{thm:cdk} into a tool is the variational
Hodge conjecture of Grothendieck: if $\cX\to S$ is smooth projective over a
connected base, $\alpha$ is a flat section of $R^{2p}f_{*}\QQ$ that is Hodge at
every point, and $\alpha_{s_{0}}$ is algebraic at one point $s_{0}$, then
$\alpha_{s}$ is algebraic at every point. This is open. What is known is the
infinitesimal statement under a semiregularity hypothesis: if the cycle at
$s_{0}$ is semiregular in the sense of Bloch \cite{Blo72}, then it deforms to
first order, and hence formally, along the locus where its class stays Hodge.
That is the theorem discussed in \Cref{sec:semireg}, and
\Cref{cor:countnotbind} records that its numerical hypothesis is not what
obstructs the present problem.
\end{remark}

\begin{remark}[Where this leaves the construction of
\texorpdfstring{\Cref{sec:construction}}{Section 14}]\label{rem:cdkconstruction}
The construction of \Cref{sec:construction} produces, for every $n$, every $d$
and every nondegenerate $\beta$, an abelian variety of Weil type together with
a sheaf whose second Chern class is computed in closed form. Its output is
algebraic on the split locus. The missing ingredient recorded in
\Cref{rem:deformation} is a statement along the algebraic locus of
\Cref{cor:weilalgebraic}: one needs the Hodge type of the deformed object to
be controlled as the parameter moves off the split locus. \Cref{thm:cdk}
guarantees that the locus one is deforming along is an algebraic variety, which
is what makes the question well posed; it does not answer it. Stating this
precisely is the purpose of \Cref{q:higher}.
\end{remark}
